\section{Проблема и постановка задачи}

\begin{frame}[plain,noframenumbering]
    \begin{center}
        \Huge Проблема и постановка задачи
    \end{center}
\end{frame}

\begin{frame}
	\frametitle{Домен исследования: AMR-системы}
	\small
	\begin{columns}[T,onlytextwidth]
		\begin{column}{0.52\textwidth}
			\begin{block}{AMR-система}
				\begin{center}
					\maybeincludegraphics[width=0.86\linewidth,height=0.32\textheight,keepaspectratio]{Dissertation/images/ch1_warehouse_robotics}
				\end{center}
				\begin{itemize}
					\item Множество автономных мобильных роботов формирует поток вычислительных задач.
					\item Задачи различаются по ресурсоёмкости, критичности, приоритету и допустимой задержке.
				\end{itemize}
			\end{block}
		\end{column}
		\begin{column}{0.44\textwidth}
			\begin{block}{Типовые задачи}
				\begin{itemize}
					\item локализация и построение карты;
					\item обработка лидаров и камер;
					\item обнаружение препятствий;
					\item перепланирование маршрута;
					\item инвентаризация и диагностика.
				\end{itemize}
				\textbf{Цикл} <<восприятие $\rightarrow$ вычисление $\rightarrow$ действие>>
			\end{block}
				\begin{block}{Проблема}
				Где выполнять неоднородный поток вычислительных задач?
			\end{block}
		\end{column}
	\end{columns}
\end{frame}

\begin{frame}
    \frametitle{Почему edge computing требуется для AMR-систем}
    \begin{center}
    	\maybeincludegraphics[width=0.95\linewidth,height=0.4\textheight,keepaspectratio]{Dissertation/images/ch1_edge_need_concept}
    \end{center}
   \small
    \begin{block}{Проблема}
    	Как распределять общий ограниченный ресурс на периферии сети?
    \end{block}
\end{frame}

\begin{frame}
	\frametitle{Правило распределения ресурсов}
	\small
	Назначение одной задачи на узел изменяет доступность ресурсов для остальных задач, то есть создаёт \textbf{внешний эффект}. 
	\begin{center}
		\maybeincludegraphics[width=0.95\linewidth,height=0.95\textheight,keepaspectratio]{Dissertation/images/ch1_edge_amr}
	\end{center}
	\begin{block}{Проблема}
		Как построить локально обоснованное правило распределение и оценить внешний эффект назначения?
	\end{block}
\end{frame}

\begin{frame}
	\frametitle{Динамичность системы}
	AMR-система с периферийными вычислениями является \textbf{последовательной динамической системой}. На каждом шаге требуется знать:
	\begin{enumerate}
		\item Какую часть ресурсов узел может предоставить новым задачам?
		\item Какую часть резервирует под критические классы?
		\item Когда переходит в режим защиты от перегрузки?
		\item Как реагирует на ухудшение связи или отказ соседних узлов?
	\end{enumerate}
	Эти решения имеют долгосрочный характер и должны вырабатываться на основе опыта взаимодействия со средой.
	\small
	\begin{block}{Проблема}
		Требуется отдельный контур адаптации. Его задача состоит не в замене механизма распределения, а в изменении условий, в которых проводится это распределение. 
	\end{block}
\end{frame}

\begin{frame}
	\frametitle{Цель и исследовательская идея}
	\small
	\textbf{Цель:} разработать модели и гибридный метод адаптивного распределения вычислительных задач автономных мобильных роботов между периферийными узлами.
	\medskip
	
	\begin{block}{Основная идея}
		Разделить две функции управления:
		\begin{enumerate}
			\item аукционный слой выбирает распределение задач и оценивает внешний эффект использования общего ресурса;
			\item обучаемый многоагентный слой меняет операционные режимы edge-узлов, адаптируя условия распределения к динамике нагрузки.
		\end{enumerate}
	\end{block}
	
	\begin{center}
		\small
		\textbf{Роботы не обучают заявки. MARL не подменяет правило распределения.}
	\end{center}
\end{frame}
\note{Ключевая формулировка: обучение не меняет смысл заявок и не заменяет аукцион. Оно управляет режимами узлов, то есть параметрами доступной ёмкости и политики допуска задач.}

\section{Предложенное решение}

\begin{frame}[plain,noframenumbering]
    \begin{center}
        \begin{center}
        	\maybeincludegraphics[width=0.8\linewidth,height=0.8\textheight,keepaspectratio]{Dissertation/images/ample_amr_logo}
        	\textit{предложенное решение}
        \end{center}
    \end{center}
\end{frame}

\begin{frame}
	\frametitle{AMPLE-AMR}
	\small
	Архитектура метода AMPLE--AMR (\textit{Adaptive Multi-agent Policy Learning with Edge Auctions for Autonomous Mobile Robots}) построена как двухуровневый контур принятия решений.
	\begin{center}
		\maybeincludegraphics[width=0.8\linewidth,height=0.8\textheight,keepaspectratio]{Dissertation/images/ch5_ample_amr_concept}
	\end{center}
\end{frame}

\begin{frame}
	\frametitle{VCG-подобный аукционный слой}
	\small
	\begin{block}{Общественное благосостояние}
		Определяется как сумма чистых вкладов назначений:
		\begin{align*}
			SW(x,t)&=\sum_{\tau\in\mathcal{T}(t)}\sum_{j=1}^{k}x_{\tau j}(t)g_{\tau j}(t) & g_{\tau j}(t)=U_{i(\tau)}(\tau,n_j,t)-C_j(\tau,r_{i(\tau)},t)
		\end{align*}
	\end{block}
	\begin{columns}[T,onlytextwidth]
		\begin{column}{0.48\textwidth}
			\begin{block}{Определение победителей}
				\begin{align*}
					x^*(t)\in\arg\max_{x\in\mathcal{X}(t,\kappa(t))}SW(x,t).
				\end{align*}
			\end{block}
			Максимизирует качество текущего распределения, но не отвечает за долговременное изменение очередей, отказов и будущей нагрузки.
		\end{column}
		\begin{column}{0.48\textwidth}
			\begin{block}{VCG-подобный платеж}
				\begin{align*}
					p_h(t)=SW^*_{-h}(t)-\sum_{\ell\in\mathcal{H},\ell\neq h}v_{\ell}(x^*(t),t).
				\end{align*}
				Рассматривается не как денежный перевод, а как внутренняя теневая цена, показывающая, какой ущерб назначение задачи создаёт для остальных задач системы.
			\end{block}
		\end{column}
	\end{columns}
\end{frame}

\begin{frame}
	\frametitle{MARL-контур}
	\small
	MARL-контур не назначает задачи напрямую. Он выбирает режимы узлов:
	\begin{align*}
		\pi_j( h^j_t)& \rightarrow \kappa_j(t) \rightarrow
		\bar{A}^{r}_j(t,\kappa_j) \rightarrow
		\mathcal{X}(t,\kappa(t)) \rightarrow x^*(t)
	\end{align*}
	
	\textbf{Режимы узлов}:
	\begin{itemize}
		\item $\kappa^{safe}$: узел ограничивает приём новых задач при высокой загрузке, риске нарушения сроков или признаках деградации;
		\item $\kappa^{normal}$: соответствует штатной работе с умеренным экспонированием свободной ёмкости;
		\item $\kappa^{aggressive}$: соответствует расширенному приёму задач при благоприятном состоянии узла и низкой очереди;
		\item $\kappa^{reserve}$: сохраняет часть ёмкости для возможных будущих критических задач;
		\item  $\kappa^{critical}$: отдаёт предпочтение задачам, связанным с безопасностью, локализацией, обходом препятствий и срочным перепланированием.
	\end{itemize}
\end{frame}

\begin{frame}
	\frametitle{AMPLE-AMR: гибридный цикл управления}
	\small
	\begin{center}
		\maybeincludegraphics[width=0.94\linewidth,height=0.4\textheight,keepaspectratio]{Dissertation/images/ch5_ample_architecture}
	\end{center}
	\begin{enumerate}
		\item Роботы формируют поток вычислительных задач.
		\item Периферийные узлы наблюдают локальное состояние и выбирают операционные режимы через QMIX-политику.
		\item Выбранные режимы определяют доступную ёмкость, правила допуска и ограничения распределения.
		\item VCG-подобный слой решает задачу выбора победителей и вычисляет внутренние теневые цены.
		\item Среда обновляет очереди, задержки, загрузку узлов и метрики качества обслуживания.
	\end{enumerate}
\end{frame}

\begin{frame}
	\frametitle{C-AMPLE-AMR: масштабируемая версия}
	\small
	\begin{columns}[T,onlytextwidth]
		\begin{column}{0.48\textwidth}
			\textbf{Проблема глобального раунда}
			\begin{itemize}
				\item с ростом числа роботов и задач увеличивается задача выбора победителей;
				\item время распределения начинает занимать значимую долю управляющего слота;
				\item глобальный аукцион становится потенциальным узким местом.
			\end{itemize}
		\end{column}
		\begin{column}{0.48\textwidth}
			\textbf{Идея C-AMPLE-AMR}
			\begin{itemize}
				\item построить граф периферийной инфраструктуры;
				\item провести локальные аукционы внутри кластеров;
				\item использовать кластеризацию только при достаточном масштабе.
			\end{itemize}
		\end{column}
	\end{columns}
\end{frame}

\begin{frame}
	\frametitle{Синтез подхода}
	\begin{center}
		\maybeincludegraphics[width=0.95\linewidth,height=0.95\textheight,keepaspectratio]{Dissertation/images/ch1_method_logic}
	\end{center}
\end{frame}

\section{Математическая формализация}

\begin{frame}[plain,noframenumbering]
	\begin{center}
		\Huge Математическая формализация
	\end{center}
\end{frame}

\begin{frame}
	\frametitle{Модели}
	\small
	\begin{center}
		\maybeincludegraphics[width=\linewidth,height=0.68\textheight,keepaspectratio]{Dissertation/images/ch2_model_mapping}
	\end{center}
\end{frame}


\begin{frame}
	\frametitle{Графовая модель инфраструктуры}
	\footnotesize
	\begin{columns}[T,onlytextwidth]
		\begin{column}{0.48\textwidth}
			Система может быть представлена в виде графа:
			\begin{align*}
				G(t)&=(\mathcal{V},a(t),w(t)) & \mathcal{V}&=\mathcal{R}\cup\mathcal{N},
			\end{align*}
			\begin{center}
				\maybeincludegraphics[width=0.94\linewidth,height=0.4\textheight,keepaspectratio]{Dissertation/images/ch2_system_graph}
			\end{center}
		\end{column}
		\begin{column}{0.48\textwidth}
			\small
			\begin{itemize}
				\item дискретное время $t=0,1,2,\ldots$;
				\item роботы $\mathcal{R}=\{r_1,\ldots,r_m\}$;
				\item периферийные узлы $\mathcal{N}=\{n_1,\ldots,n_k\}$;
				\item индикатор достижимости $a_{uv}(t)$;
				\item сетевые задержки $w_{uv}(t)$.
			\end{itemize}
		\end{column}
	\end{columns}
\end{frame}

\begin{frame}
	\frametitle{Модель ресурсов и задач}
	\footnotesize
	\begin{columns}[T,onlytextwidth]
		\begin{column}{0.52\textwidth}
			\textbf{Поток задач}
			\begin{align*}
				\mathcal{T}(t)&=\bigcup_{i=1}^{m}\mathcal{T}_i(t) &
				\mathcal{T}_i(t)&=\{\tau_{i1}(t),\tau_{i2}(t),\ldots\}
			\end{align*}
			\begin{itemize}
				\item $c_{\tau}$ --- класс задачи,
				\item $b^{in}_{\tau}$ и $b^{out}_{\tau}$ --- объёмы входных и выходных данных,
				\item $q^{cpu}_{\tau}$, $q^{mem}_{\tau}$ и $q^{gpu}_{\tau}$ --- требования к вычислительным ресурсам,
				\item $\Delta_{\tau}$ --- допустимый срок обработки,
				\item $\pi_{\tau}$ --- приоритет.
			\end{itemize}
			\textbf{Характеристики узла}
			\begin{itemize}
				\item $A_j(t) = (A^{cpu}_j(t),A^{mem}_j(t),A^{gpu}_j(t)) $ --- ёмкость узла,
				\item $Q_j(t)$ --- очередь обслуживания;
				\item $f_j(t)$ --- индикатор отказа;
				\item $x_{\tau j}(t)$ --- переменная назначения;
			\end{itemize}
		\end{column}
		\begin{column}{0.44\textwidth}
			\textbf{Допустимое распределение}  $\mathcal{X}(t)$
			
			\textit{Единственность назначения}:
			\begin{align*}
				\sum_{j=1}^{k}x_{\tau j}(t)&\leq 1, && \forall \tau\in\mathcal{T}(t)
			\end{align*}
			
			\textit{Ëмкость узла}:
			\begin{align*}
				\sum_{\tau\in\mathcal{T}(t)}x_{\tau j}(t)q^{r}_{\tau}&\leq A^{r}_j(t,\kappa_j), && \forall j
			\end{align*}
			
			\textit{Сетевая достижимость}:
			\begin{align*}
				x_{\tau j}(t)&\leq a_{i(\tau)j}(t)
			\end{align*}
			
			\textit{Отказ узла}:
			\begin{align*}
				x_{\tau j}(t)\leq 1-f_j(t)
			\end{align*}
			
			\textit{Срок выполнения}:
			\begin{align*}
				x_{\tau j}(t)T_{\tau j}(t)&\leq \Delta_{\tau}, && \forall \tau:c_{\tau}\in\mathcal{C}^{hard}
			\end{align*}
		\end{column}
	\end{columns}
\end{frame}

\begin{frame}
	\frametitle{Модель задержки выполнения задачи}
	\small
	\begin{align*}
		T_{\tau j}(t)=T^{net}_{ij,\tau}(t)+T^{queue}_{j}(t)+T^{proc}_{\tau j}(t),
	\end{align*}
	
	Сетевая составляющая
	\begin{align*}
		T^{net}_{ij,\tau}(t)=d_{ij}(t)+\frac{b^{in}_{\tau}}{B^{up}_{ij}(t)+\varepsilon}+\frac{b^{out}_{\tau}}{B^{down}_{ji}(t)+\varepsilon}
	\end{align*}
	
	Время ожидания в очереди
	\begin{align*}
		T^{queue}_{j}(t)=\frac{\sum_{\tau'\in Q_j(t)} q^{cpu}_{\tau'}}{F^{cpu}_{j}(t)+\varepsilon},
	\end{align*}
	
	Время непосредственной обработки
	\begin{align*}
		T^{proc}_{\tau j}(t)=\frac{q^{cpu}_{\tau}}{F^{cpu}_{j}(t)+\varepsilon}+\mathbb{I}\{q^{gpu}_{\tau}>0\}\frac{q^{gpu}_{\tau}}{F^{gpu}_{j}(t)+\varepsilon}.
	\end{align*}
\end{frame}

\begin{frame}
	\frametitle{Модель полезности назначения}
	\small
	\begin{align*}
		U_i(\tau,n_j,t) &= V_i(\tau)\psi_{c_{\tau}}(T_{\tau j}(t),\Delta_{\tau})-E^{tx}_{i\tau j}(t)
	\end{align*}
	
	Базовая ценность задачи для робота-источника
	\begin{align*}
		V_i(\tau) &= \beta_i v_{c_{\tau}}\pi_{\tau}\left(1+\theta_{c_{\tau}}\log\left(1+\frac{W_{\tau}}{W_0}\right)\right),\\[-1mm]
		W_{\tau} &= \omega_{cpu}q^{cpu}_{\tau}+\omega_{mem}q^{mem}_{\tau}+\omega_{gpu}q^{gpu}_{\tau}+\omega_{data}b^{in}_{\tau}
	\end{align*}
	
	Функция деградации полезности при задержке
	\begin{align*}
		\psi_{c_{\tau}}(T_{\tau j}(t),\Delta_{\tau}) &=
		\exp\left(-\lambda_{c_{\tau}}\frac{T_{\tau j}(t)}{\Delta_{\tau}+\varepsilon}\right)
	\end{align*}
	
	Энергетические затраты передачи
	\begin{align*}
		E^{tx}_{i\tau j}(t)=e^{tx}_i b^{in}_{\tau} T^{net}_{ij,\tau}(t)
	\end{align*}
\end{frame}
\note{Назначение оценивается не только по вычислительной близости узла, а по системной полезности с учётом задержки, энергии передачи и риска нарушения срока.}

\begin{frame}
	\frametitle{Модель стоимости узла}
	\small
	\begin{align*}
		C_j(\tau,r_i,t)&=C^{comp}_j(\tau)+C^{net}_j(\tau,r_i,t)+C^{over}_j(\tau,t)+C^{rel}_j(t)
	\end{align*}
	
	Вычислительная стоимость
	\begin{align*}
		C^{comp}_j(\tau)&=\rho^{cpu}_j q^{cpu}_{\tau}+\rho^{mem}_j q^{mem}_{\tau}+\rho^{gpu}_j q^{gpu}_{\tau}
	\end{align*}
	
	Сетевая стоимость
	\begin{align*}
		C^{net}_j(\tau,r_i,t)&=\rho^{net}_j\left(b^{in}_{\tau}+b^{out}_{\tau}\right)w_{ij}(t)
	\end{align*}
	
	Стоимость перегрузки
	\begin{align*}
		C^{over}_j(\tau,t)&=\mu^{cpu}_j\left(\frac{L^{cpu}_j(t)+q^{cpu}_{\tau}}{C^{cpu}_j+\varepsilon}\right)^2+
		\mu^{mem}_j\left(\frac{L^{mem}_j(t)+q^{mem}_{\tau}}{C^{mem}_j+\varepsilon}\right)^2+{}\\[-1mm]
		&\quad{}+\mu^{gpu}_j\left(\frac{L^{gpu}_j(t)+q^{gpu}_{\tau}}{C^{gpu}_j+\varepsilon}\right)^2
	\end{align*}
	
	Компонент надёжности
	\begin{align*}
		C^{rel}_j(t)&=\rho^{fail}_j f_j(t)
	\end{align*}
\end{frame}

\begin{frame}
	\frametitle{Дизайн экономических механизмов}
	\small
	\textbf{Механизм} -- задает критерий выбора исхода, а также правило, позволяющее количественно оценить влияние участника на остальных:
	\begin{align*}
		M &= (g,p) &
		g &: \Theta_1 \times \ldots \times \Theta_{|\mathcal{H}|} \rightarrow X & p_h : \Theta_1 \times \ldots \times \Theta_{|\mathcal{H}|} \rightarrow \mathbb{R}
	\end{align*}
	где $\mathcal{H}$ -- множество участников механизма, $X$ -- множество допустимых исходов (назначение задач на узлы), $g$ -- правило выбора исхода, а $p_h$ -- платёж или внутренняя компенсация участника $h$.
	
	\textbf{Тип участника} $h \in \mathcal{H}$ задаёт его локальную информацию $\theta_h$:
	\begin{itemize}
		\item \textit{Для робота} -- класс задачи, ценность результата, срок выполнения и требования к ресурсам;
		\item \textit{Для периферийного узла} -- доступная ёмкость, очередь, стоимость обслуживания и состояние каналов связи.
	\end{itemize}
	
	\textbf{Удовлетворенность участника}
	\begin{align*}
		u_h(\theta_h,\widehat{\theta}) = v_h(g(\widehat{\theta}),\theta_h) - p_h(\widehat{\theta})
	\end{align*}
	
	\textbf{Общественное благосостояние}:
	\begin{align*}
		SW(x) &= \sum_{h \in \mathcal{H}} v_h(x) & v_{r_i}(x) &= U_i(x) & v_{n_j}(x) &= - C_j(x)
	\end{align*}
\end{frame}
\note{Вместо анализа поведения участников при заданных правилах оно строит такие правила взаимодействия, при которых локальные сообщения участников приводят к требуемому системному результату. Для настоящей работы это означает задание множества допустимых распределений задач, формата заявок, правила выбора распределения и правила вычисления внутренних компенсаций.}

\begin{frame}
	\frametitle{VCG-подобный аукционный слой}
	\small
	\begin{block}{Общественное благосостояние}
		Определяется как сумма чистых вкладов назначений:
		\begin{align*}
			SW(x,t)&=\sum_{\tau\in\mathcal{T}(t)}\sum_{j=1}^{k}x_{\tau j}(t)g_{\tau j}(t) & g_{\tau j}(t)=U_{i(\tau)}(\tau,n_j,t)-C_j(\tau,r_{i(\tau)},t)
		\end{align*}
	\end{block}
	\begin{columns}[T,onlytextwidth]
		\begin{column}{0.48\textwidth}
			\begin{block}{Определение победителей}
				\begin{align*}
					x^*(t)\in\arg\max_{x\in\mathcal{X}(t,\kappa(t))}SW(x,t).
				\end{align*}
			\end{block}
			Максимизирует качество текущего распределения, но не отвечает за долговременное изменение очередей, отказов и будущей нагрузки.
		\end{column}
		\begin{column}{0.48\textwidth}
			\begin{block}{VCG-подобный платеж}
				\begin{align*}
					p_h(t)=SW^*_{-h}(t)-\sum_{\ell\in\mathcal{H},\ell\neq h}v_{\ell}(x^*(t),t).
				\end{align*}
				Рассматривается не как денежный перевод, а как внутренняя теневая цена, показывающая, какой ущерб назначение задачи создаёт для остальных задач системы.
			\end{block}
		\end{column}
	\end{columns}
\end{frame}

\begin{frame}
	\frametitle{MARL-контур}
	\small
	\begin{center}
		\maybeincludegraphics[width=\linewidth,height=0.68\textheight,keepaspectratio]{Dissertation/images/ch4_adaptive_control}
	\end{center}
\end{frame}

\begin{frame}
	\frametitle{MARL-контур: накопление опыта}
	\small
	\begin{columns}[T,onlytextwidth]
		\begin{column}{0.52\textwidth}
			\begin{block}{Глобальное состояние $s_t$}
				\begin{itemize}
					\item $G(t)$ --- граф текущей инфраструктуры и сетевой достижимости;
					\item $L_j(t)$ --- вектор загрузки периферийного узла $n_j$ по видам ресурсов;
					\item $Q_j(t)$ --- очередь задач узла $n_j$;
					\item $F_j(t)$ --- состояние исправности, деградации или отказа узла;
					\item $\mathcal{T}(t)$ --- множество задач, поступивших в текущем раунде;
					\item $d_{ij}(t)$ --- оценка сетевой задержки между роботом $r_i$ и узлом $n_j$;
					\item $Z_t$ --- агрегированные признаки предыдущих раундов.
				\end{itemize}
			\end{block}
		\end{column}
		\begin{column}{0.44\textwidth}
			\begin{block}{Локальные наблюдения}
				\begin{align*}
					o^j_t &= \Omega_j(s_t) = \left(o^{j,tech}_t,g^j_t\right),
				\end{align*}
				\begin{itemize}
					\item $o^{j,tech}_t$ содержит технические признаки узла;
					\item $g^j_t$ --- агрегированные аукционные признаки предыдущего распределения.
				\end{itemize}
			\end{block}
			\begin{block}{Локальная история}
				\begin{align*}
					h^j_t &= \left(o^j_0,\kappa_j(0),r_0, o^j_1,\kappa_j(1),r_1,\ldots,o^j_t\right)
				\end{align*}
			\end{block}
		\end{column}
	\end{columns}
\end{frame}

\begin{frame}
	\frametitle{MARL-контур: операционные режимы}
	\small
	MARL-контур не назначает задачи напрямую. Он выбирает режимы узлов:
	\begin{align*}
		\pi_j( h^j_t)& \rightarrow \kappa_j(t) \rightarrow
		\bar{A}^{r}_j(t,\kappa_j) \rightarrow
		\mathcal{X}(t,\kappa(t)) \rightarrow x^*(t)
	\end{align*}
	
	\textbf{Режимы узлов}:
	\begin{itemize}
		\item $\kappa^{safe}$: узел ограничивает приём новых задач при высокой загрузке, риске нарушения сроков или признаках деградации;
		\item $\kappa^{normal}$: соответствует штатной работе с умеренным экспонированием свободной ёмкости;
		\item $\kappa^{aggressive}$: соответствует расширенному приёму задач при благоприятном состоянии узла и низкой очереди;
		\item $\kappa^{reserve}$: сохраняет часть ёмкости для возможных будущих критических задач;
		\item  $\kappa^{critical}$: отдаёт предпочтение задачам, связанным с безопасностью, локализацией, обходом препятствий и срочным перепланированием.
	\end{itemize}
\end{frame}

\begin{frame}
	\frametitle{MARL-контур: командная цель}
	\small
	Совместная политика периферийных узлов
	\begin{align*}
		\pi &= \left(\pi_1,\ldots,\pi_k\right)
	\end{align*}
	
	Цель обучения совместной политики периферийных узлов:
	\[
	J(\pi)=\mathbb{E}_{\pi}\left[\sum_{t=0}^{\infty}\gamma^t r_t\right],
	\qquad
	\pi^*\in\operatorname*{arg\,max}_{\pi}J(\pi).
	\]
	Базовый командный сигнал:
	\begin{align*}
		r_t ={}& \alpha_{sw}\widetilde{SW}(t)
		-\alpha_{viol}\mathrm{Viol}(t)
		-\alpha_{drop}\mathrm{Drop}(t)
		-\alpha_{wait}\mathrm{Wait}(t)
	\end{align*}
	С учётом аукционных признаков:
	\[
	\tilde{r}_t=(1-\lambda_{auc})r_t+\lambda_{auc}r^{auc}_t,
	\quad \lambda_{auc}\in[0,1].
	\]
\end{frame}

\begin{frame}
	\frametitle{MARL-контур: QMIX-формализация}
	\small
	Для каждого узла $n_j$ локальная функция ценности оценивает режим $\kappa_j$ по локальной истории $h^j_t$:
	\begin{align*}
		Q_j(h^j_t,\kappa_j;\theta_j),\qquad
		Q_{tot}(s_t,\kappa_t)=f_{mix}(Q_1,\ldots,Q_k,s_t;\theta_{mix}).
	\end{align*}
	Ограничение монотонности QMIX:
	\[
	\frac{\partial Q_{tot}}{\partial Q_j}\geq 0,\quad \forall j.
	\]
	Тогда локальный жадный выбор согласован с совместной функцией ценности:
	\[
	\kappa_j^*(t)=\operatorname*{arg\,max}_{\kappa_j\in\mathcal{K}_j}Q_j(h^j_t,\kappa_j;\theta_j).
	\]
	Обучение по TD-ошибке:
	\begin{align*}
		y_t&=\tilde{r}_t+\gamma(1-done_t)
		\max_{\kappa'_{t+1}}Q^{-}_{tot}(s_{t+1},\kappa'_{t+1}),\\[-1mm]
		\mathcal{L}(\theta)&=
		\mathbb{E}_{e_t\sim\mathcal{B}}
		\left[\left(y_t-Q_{tot}(s_t,\kappa_t;\theta)\right)^2\right].
	\end{align*}
\end{frame}

\begin{frame}
	\frametitle{MARL-контур: QMIX-сети}
	\small
	\begin{center}
		\maybeincludegraphics[width=\linewidth,height=0.8\textheight,keepaspectratio]{Dissertation/images/ch4_arch_qmix}
	\end{center}
\end{frame}

\section{Прототип и экспериментальная проверка}

\begin{frame}[plain,noframenumbering]
    \begin{center}
        \Huge Прототип и экспериментальная проверка
    \end{center}
\end{frame}

\begin{frame}
    \frametitle{Программный прототип}
    \small
    \begin{columns}[T,onlytextwidth]
        \begin{column}{0.48\textwidth}
            \textbf{Реализованные компоненты}
            \begin{itemize}
                \item генератор складских сценариев;
                \item модель AMR-задач и edge-узлов;
                \item VCG-подобный распределитель;
                \item QMIX-контур обучения режимов;
                \item кластерный режим C-AMPLE-AMR;
                \item экспорт CSV, графиков и \LaTeX-таблиц.
            \end{itemize}
        \end{column}
        \begin{column}{0.48\textwidth}
            \textbf{Сравниваемые методы}
            \begin{enumerate}
                \item Fixed-Heuristic;
                \item Fixed-VCG;
                \item QMIX-Heuristic;
                \item AMPLE-AMR;
                \item C-AMPLE-AMR.
            \end{enumerate}
            \medskip
            \begin{block}{Принцип сравнения}
                Все методы проверяются в одинаковой среде, на одинаковых сценариях и при одинаковых начальных значениях генератора случайных чисел.
            \end{block}
        \end{column}
    \end{columns}
\end{frame}

\begin{frame}
	\frametitle{Складская среда}
	\small
	\begin{center}
		\maybeincludegraphics[width=\linewidth,height=0.68\textheight,keepaspectratio]{Dissertation/images/ch6_warehouse_scenario}
	\end{center}
\end{frame}

\begin{frame}
	\frametitle{Проверка экспериментальных гипотез}
	\scriptsize
	\begin{tabular}{p{0.12\linewidth}p{0.49\linewidth}p{0.28\linewidth}}
		\toprule
		\textbf{Гипотеза} & \textbf{Смысл} & \textbf{Итог} \\
		\midrule
		H1 & эффективность гибридизации & частично подтверждена \\
		H2 & польза обучаемого выбора режимов & подтверждена \\
		H3 & польза VCG-подобного распределения & не подтверждена в сильной форме \\
		H4 & устойчивость к пиковой нагрузке & подтверждена \\
		H5 & устойчивость к неоднородности ресурсов & частично подтверждена \\
		H6 & устойчивость к деградации связи и отказам & подтверждена \\
		H7 & масштабируемость C-AMPLE-AMR & условно подтверждена \\
		\bottomrule
	\end{tabular}
	\medskip
	
	\begin{block}{Интерпретация}
		Главный практический эффект текущего прототипа связан с обучаемым выбором режимов edge-узлов. VCG-подобный слой даёт интерпретируемость, внутренние цены и формальную структуру распределения, но не всегда превосходит сильную жадную эвристику по метрикам качества.
	\end{block}
\end{frame}

\begin{frame}
	\frametitle{Результаты валидации}
	\small
	\begin{itemize}
		\item Обучение политики не вырождается в фиксированный выбор операционного режима; 
		\item AMPLE-AMR выигрывает по общественному благосостоянию; 
		\item C-AMPLE-AMR нужен только на крупном масштабе;
		\item VCG-слой полезен прежде всего как интерпретируемая структура и внутренние цены, но не как гарантированное улучшение над сильной эвристикой в каждом сценарии
	\end{itemize}
\end{frame}
